% !TEX root = ../Franciszek_Zabierowski_mgr.tex
%autorzy szablonu: Przemysław Michalski, z inspiracji Agaty Dorau, przemyslaw.michalski@pw.edu.pl

\linespread{1.3} %interlinia 1.5

\usepackage[MeX]{polski}

%fonty
\usepackage[T1]{fontenc}
\usepackage[utf8]{inputenc}
\usepackage[sfdefault,scaled=.95]{FiraSans}
% \usepackage{newtxsf}

%definicje kolorów
\usepackage[table]{xcolor}
\definecolor{sliwka}{RGB}{150,95,119} %kolor śliwkowy, zgodny z zaleceniami
\definecolor{grafit}{RGB}{60,60,60} %kolor grafitowy, zgodny z zaleceniami (przy druku monochromatycznym)

%linki w PDFie - linkowane są przypisy oraz bibliografia
\usepackage{hyperref}
\hypersetup{
	pdftitle = {}, %tytuł dokumentu PDF
	pdfauthor = {}, %autor PDFa
	colorlinks,
	linktoc=page,
	linkcolor=sliwka,
	citecolor=sliwka,
	filecolor=sliwka,
	urlcolor=sliwka
}

\usepackage{amsmath,amsfonts,amssymb} %fonty matematyczne
\usepackage{braket}
\usepackage{icomma} %brak spacji po przecinku w trybie matematycznym (zwykle przecinek jest traktowany jako rozdział między kolejnymi trzema cyframi)
\usepackage{textcomp} %the package supports the Text Companion fonts, which provide many text symbols
\usepackage{gensymb} %stopien Celsjusza (\celcius), stopien (\degree), mikro (\micro), om (\textohm) i promil (\perthousand) w trybie tekstowym

\usepackage{graphicx}

%marginesy
\usepackage[left=3cm,right=2cm,top=2.5cm,bottom=2.5cm]{geometry} 

%nagłówek i stopka
\usepackage{fancyhdr}
\pagestyle{fancy}
\renewcommand{\headrulewidth}{0pt} %brak linii w nagłówku
\renewcommand{\footrulewidth}{1.5pt} %linia w stopce
\renewcommand{\footrule}{{\color{sliwka}%
	\vskip-\footruleskip\vskip-\footrulewidth
	\hrule width\headwidth height\footrulewidth\vskip\footruleskip}} %kolor linii w stopce

\fancyhead{} %pusty nagłówek
\cfoot{} %pusta środkowa część stopki
\fancyfoot[LE,RO]{\thepage} %numer strony w stopce na parzystej stronie z lewej, na nieparzystej - z prawej

%listy - kształt punktów
\renewcommand\labelitemi{\color{sliwka}{\tiny$\blacksquare$}} %1. poziom
\renewcommand\labelitemii{\color{sliwka}{\tiny$\square$}} %2. poziom
\renewcommand\labelitemiii{\color{sliwka}{\tiny$\bullet$}} %3. poziom

%wcięcia w listach
\usepackage{enumitem}
\setlist[itemize]{noitemsep, topsep=0pt, leftmargin=1cm, itemindent=0cm}

%stylistyka
\usepackage{nowidow}
\widowpenalty=10000 %nie pozostawia wdów na końcu strony
\clubpenalty=10000 %nie pozostawia sierot
\brokenpenalty 10000 %nie dzieli stron, gdy przenoszony jest wyraz
\sloppy % zakaz wydłużania linii poza margines
\frenchspacing %,,kontynentalne" zwyczaje typograficzne

%odstępy
\usepackage[skip=4pt,indent=0cm,parfill=0pt]{parskip}

%obrazki/tabele/równania
\usepackage{array}
\usepackage{multicol}
\usepackage{multirow}
\renewcommand*{\figurename}{Rys.} 
\renewcommand*{\tablename}{Tab.}
\setlength{\arrayrulewidth}{1pt} %grubość linii tabeli
\arrayrulecolor{sliwka} %kolor linii tabeli
\renewcommand\thefootnote{\textcolor{sliwka}{\arabic{footnote}}} %kolor numeru przypisu
\usepackage[font={footnotesize},justification=centerlast]{caption} %czcionka 9 pkt w podpisach tabeli/obrazków
\numberwithin{equation}{chapter} %numerowanie równań od numeru rozdziału
\newcolumntype{M}[1]{>{\centering\arraybackslash}m{#1}}
\usepackage{subfig}
\usepackage{subcaption}
\usepackage{wrapfig}
\usepackage{lipsum} %generuje tekst

\let\lll=\relax

% ---------- DUŻE SUMY I CAŁKI ----------
\usepackage{relsize}
% Użycie: $\bigsum_{n=0}^{\infty}$, $\bigint_0^1 f(x)\,dx$, itp.
\newcommand{\LargeSum}[1][1.5]{\mathop{\vcenter{\hbox{\scalebox{#1}{$\displaystyle\sum$}}}}\displaylimits}
\newcommand{\LargeInt}[1][1.25]{\mathop{\vcenter{\hbox{\scalebox{#1}{$\displaystyle\int$}}}}\displaylimits}
\newcommand{\LargeIInt}[1][1.25]{\mathop{\vcenter{\hbox{\scalebox{#1}{$\displaystyle\iint$}}}}\displaylimits}

\usepackage{comment}
\usepackage{float}
\usepackage{verbatim}


\captionsetup[sub]{justification=centering,list=yes}
\newtheorem{theorem}{Twierdzenie}

\usepackage{graphicx} % \scalebox

% Dostosowanie etykiet dla polskiego
\usepackage{cleveref}
\crefname{chapter}{rozdział}{rozdziały}
\crefname{section}{sekcja}{sekcje}
\crefname{subsection}{podsekcja}{podsekcje}
\crefname{equation}{równanie}{równania}
\crefname{figure}{rysunek}{rysunki}
\crefname{table}{tabela}{tabele}


\newcommand{\forceoddpage}{%
	\clearpage
	\ifodd\value{page}\else
		\thispagestyle{empty}\null\newpage
	\fi
}

\pagecolor{black}
\color{white}